\labsection{6}{Поиск экстремума функции двух переменных с помощью генетического алгоритма с бинарным представлением особей}

\labpart{Цель работы}
Модифицировать простой генетический алгоритм для многомерного поиска экстремума; изучить бинарное
кодирование нескольких переменных, операторы многоточечного и однородного кроссинговера, оператор
инверсии, методы селекции ранжированием и турниром и стратегии сокращения промежуточной популяции;
познакомиться с теоремой схем.

\labpart{Теоретические сведения}
Подразделы~\ref{sec:ga}.5 и~\ref{sec:ga}.7, раздел~\ref{sec:operators} теоретической части: кодирование
нескольких переменных и предел точности сетки, теорема схем (\ref{eq:schema}), селекция ранжированием
(\ref{eq:rank}) и турниром, многоточечный и однородный кроссинговер, инверсия, стратегии сокращения
популяции.

\labpart{Постановка задачи}
Дана функция $y = f(x_1, x_2)$ (таблица~\ref{tab:l6-var})~--- одна из стандартных тестовых функций
численной оптимизации. Найти её минимум при $x_i \in [a; b]$ генетическим алгоритмом с бинарным
представлением особей и графически проиллюстрировать динамику поиска точки экстремума на плоскости
$X_1 X_2$. В качестве оператора кроссинговера использовать многоточечный или однородный кроссинговер,
в качестве оператора мутации~--- инверсию, селекцию выполнять ранжированием или турниром. Провести
эксперименты с разными методами сокращения промежуточной популяции.

\begin{table}[!htb]
\caption{Варианты заданий к лабораторным работам №6 и~№7 (часть А)}
\label{tab:l6-var}
\centering\small
\begin{tabularx}{\textwidth}{c>{\raggedright}p{3.1cm}Xc}
\toprule
№ & Функция & Формула & $x_i$ \\
\midrule
1 & 2-я де Йонга (Розенброка) & $100(x_1^2 - x_2)^2 + (1 - x_1)^2$ & $[-2,048; 2,048]$ \\
2 & 6-я Шаффера & $0,5 + \dfrac{\sin^2\sqrt{x_1^2 + x_2^2} - 0,5}{\bigl(1 + 0,001(x_1^2 + x_2^2)\bigr)^2}$ & $[-100; 100]$ \\
3 & 7-я Шаффера & $(x_1^2 + x_2^2)^{0,25}\bigl(\sin^2\bigl(50(x_1^2 + x_2^2)^{0,1}\bigr) + 1\bigr)$ & $[-100; 100]$ \\
4 & Голдстейна~--- Прайса & $\bigl(1 + (x_1 + x_2 + 1)^2(19 - 14x_1 + 3x_1^2 - 14x_2 + 6x_1x_2 + 3x_2^2)\bigr)\times$ \newline
  $\times\bigl(30 + (2x_1 - 3x_2)^2(18 - 32x_1 + 12x_1^2 + 48x_2 - 36x_1x_2 + 27x_2^2)\bigr)$ & $[-2; 2]$ \\
5 & Изома & $-\cos x_1 \cos x_2\, e^{-(x_1 - \pi)^2 - (x_2 - \pi)^2}$ & $[-100; 100]$ \\
6 & 1-я Бохачевского & $x_1^2 + 2x_2^2 - 0,3\cos(3\pi x_1) - 0,4\cos(4\pi x_2) + 0,7$ & $[-50; 50]$ \\
7 & 2-я Бохачевского & $x_1^2 + 2x_2^2 - 0,3\cos(3\pi x_1)\cos(4\pi x_2) + 0,3$ & $[-50; 50]$ \\
8 & Растригина & $20 + \sum_{i=1}^{2}\bigl(x_i^2 - 10\cos(2\pi x_i)\bigr)$ & $[-5,12; 5,12]$ \\
\bottomrule
\end{tabularx}
\end{table}

Функции подобраны так, чтобы проверить разные трудности поиска (рисунок~\ref{fig:l6-functions}):
узкий изогнутый овраг (1), множество концентрических колец локальных минимумов (2, 3), резкие перепады
значений (4), <<иголка>> глобального минимума на плоском фоне (5), регулярная решётка локальных минимумов
(6--8). Глобальные минимумы: 0 в точке $(1; 1)$ для функции~1, 3 в точке $(0; -1)$ для функции~4, $-1$ в
точке $(\pi; \pi)$ для функции~5 и 0 в начале координат для остальных.

\begin{figure}[!htb]
\centering
\includegraphics[width=0.88\textwidth]{\figdir/test_functions.png}
\caption{Линии уровня тестовых функций (темнее~--- меньше; для функций 5--7 показан фрагмент области
вблизи минимума)}
\label{fig:l6-functions}
\end{figure}

\labpart{Требования к программе}
Исходные данные: размер популяции, длина кода одной переменной, вероятности кроссинговера и мутации,
максимальное число поколений, метод селекции и стратегия сокращения популяции. В процессе поиска
отображаются: линии уровня функции на плоскости $X_1X_2$ и положение особей текущей популяции; графики
лучшего и среднего значений функции по поколениям; координаты лучшей особи и значение функции в ней.

\labpart{Порядок выполнения работы}
\begin{steps}
\item Построить линии уровня функции (\code{plt.contourf}), найти эталонный минимум
(\code{scipy.optimize.differential\_evolution}) и предел точности кодирования~--- минимальное значение
функции в узлах двоичной сетки вблизи глобального минимума.
\item Модифицировать программу ЛР~№5 (операторы на NumPy~--- подраздел~\ref{sec:ga-numpy}): хромосома~---
конкатенация кодов двух переменных; селекция~---
ранжирование или турнир; многоточечный или однородный кроссинговер; инверсия (совместно с побитовой
мутацией); сокращение популяции~--- элитарная схема.
\item Выполнить основной запуск с визуализацией динамики популяции.
\item Сравнить операторы мутации: инверсия с побитовой мутацией, только инверсия, только побитовая
мутация, без мутации (не менее 20 запусков каждой конфигурации).
\item Сравнить стратегии сокращения популяции: элитарную ($\mu + \lambda$), чистую замену, случайную
замену и другие по выбору.
\item Исследовать влияние размера популяции, вероятностей операторов и размера турнира.
\end{steps}

\labpart{Пример выполнения}
Вариант~3: минимум 7-й функции Шаффера на $[-100; 100]^2$, глобальный минимум $f(0; 0) = 0$. Минимумы
функции образуют концентрические кольца, сгущающиеся к началу координат, поэтому градиентный спуск
застревает на первом же кольце. При $L = 20$ бит на переменную шаг сетки $\Delta = 1,9 \cdot 10^{-4}$, но
точка 0 не представима: ближайшие узлы дают $x_i = \pm\Delta/2$, и минимальное значение функции на сетке
равно $0,0200$. Меньшего значения не получит никакой алгоритм с таким кодированием, поэтому успехом
считается $f < 0,1$, а дополнительно фиксируется доля запусков, достигших предела $0,0200$.

ГА с параметрами $N = 80$, турнир $k = 3$, однородный кроссинговер $P_c = 0,85$, инверсия
$P_{inv} = 0,2$ совместно с побитовой мутацией $1/2L$, элитарное сокращение нашёл точку
$(-9,5 \cdot 10^{-5}; 9,5 \cdot 10^{-5})$ с $f = 0,0200$~--- точно на пределе сетки
(24\,080 вычислений функции). За 10 поколений популяция стягивается с области $\pm 100$ к кольцам
радиусом несколько единиц (рисунок~\ref{fig:l6-pop}).

\begin{figure}[!htb]
\centering
\includegraphics[width=\textwidth]{\labfig{6}{03_population.png}}
\caption{Популяция ГА на линиях уровня функции Шаффера F7: поколения 0, 10 и 300 (масштаб осей уменьшается)}
\label{fig:l6-pop}
\end{figure}

Сравнение операторов мутации и стратегий сокращения (20 запусков) приведено в
таблице~\ref{tab:l6-exp}. Инверсия без побитовой мутации слаба: она лишь переставляет биты, не меняя
их состава, и не создаёт недостающих значений. Без мутации популяция вырождается. Стратегия сокращения
определяет скорость поиска: без элитизма лучшие особи теряются, и кривая лучшей особи поколения
колеблется. Рост популяции с 20 до 160 ускоряет поиск с 58 до 22 поколений, увеличение турнира с 2 до 8
особей~--- с 31 до 20 поколений.

\begin{table}[!htb]
\caption{Операторы мутации и стратегии сокращения популяции (20 запусков; поколений~--- медиана до
$f < 0,1$)}
\label{tab:l6-exp}
\centering\small
\begin{tabular}{lcccc}
\toprule
Конфигурация & Медиана $f$ & Успех, \% & До предела, \% & Поколений \\
\midrule
инверсия + побитовая, элитарное & 0,0200 & 100 & 90 & 26 \\
только инверсия & 0,0638 & 70 & 10 & 47 \\
только побитовая & 0,0200 & 100 & 80 & 29 \\
без мутации & 1,0128 & 5 & 0 & 11 \\
\midrule
чистая замена & 0,0200 & 100 & 65 & 44 \\
случайная замена & 0,0208 & 100 & 35 & 84 \\
\bottomrule
\end{tabular}
\end{table}

\labpart{Контрольные вопросы}
\begin{questions}
\item Поясните понятие схемы, её порядка и определяющей длины. Сформулируйте фундаментальную теорему
ГА и объясните вклад каждого оператора.
\item Что такое строительные блоки и неявный параллелизм?
\item Как кодируется точка с несколькими координатами? Почему ГА с двоичным кодом не может найти
точный минимум функции Шаффера на симметричной области?
\item Опишите селекцию ранжированием, турниром, усечением и локальный отбор. Как регулируется давление
отбора?
\item Опишите способы выбора пар: панмиксию, селективный выбор, инбридинг и аутбридинг.
\item Опишите многоточечный и однородный кроссинговер. Чем однородный кроссинговер отличается по
действию на схемы?
\item Почему инверсия без побитовой мутации не обеспечивает сходимости к минимуму?
\item Опишите стратегии сокращения промежуточной популяции: чистую замену, элитарную схему, случайную
замену, пропорциональную редукцию, отбор вытеснением.
\item Что такое преждевременная сходимость и как её предотвратить?
\end{questions}
